\section*{Chapitre \ref{ch:rc:counterparty-exposure} --- Exposition de contrepartie}

\begin{solution}{exo:rc:counterparty-exposure:1}
La valeur d'un swap en $t$ vaut à peu près le mouvement du taux jusqu'à $t$ multiplié par l'annuité restante~: le premier
facteur croît comme $\sqrt t$, le second se réduit à zéro, si bien que le produit culmine entre les deux (environ un
tiers du parcours pour un swap). Un swap de devises avec échange final conserve tout le risque de change du notionnel
jusqu'à l'échéance, si bien que son exposition continue de croître en $\sqrt t$.
\end{solution}

\begin{solution}{exo:rc:counterparty-exposure:2}
Avec compensation, $\max(10-7,0) = 3$~; sans compensation, la banque perd le $+10$ et doit quand même payer le
$7$ qu'elle doit~: exposition 10.
\end{solution}

\begin{solution}{exo:rc:counterparty-exposure:3}
Le collatéral requis vaut $14.3-10 = 4.3$ millions~; la banque en détient 4.0, si bien que la variation
de 0.3 million est inférieure au MTA de 0.5~: pas d'appel de marge.
\end{solution}

\begin{solution}{exo:rc:counterparty-exposure:4}
Avec un seuil nul, l'exposition collatéralisée est le mouvement de la valeur sur la période de marge
en risque, dont l'écart type croît comme la racine carrée de la période pour une diffusion~;
l'EE est proportionnelle à l'écart type. La simulation donne 1.41.
\end{solution}

\begin{solution}{exo:rc:counterparty-exposure:5}
La contrepartie s'attend à recevoir plus de la banque qu'elle ne lui doit~: son exposition à la banque
est supérieure à celle de la banque envers elle. La DVA de la banque, soit la valeur de son propre défaut
pour ses créanciers, se calcule à partir de l'ENE et est en conséquence élevée.
\end{solution}

\begin{solution}{exo:rc:counterparty-exposure:6}
Exemples spécifiques~: acheter une protection sur les banques d'un pays auprès d'une de ces banques~; un put sur
les actions d'une entreprise acheté à cette entreprise. Cas de corrélation favorable~: un producteur de matières premières qui vend sa
production à terme à la banque lui doit le plus lorsque les prix sont élevés, c'est-à-dire quand le
producteur est le plus solide.
\end{solution}

\begin{solution}{exo:rc:counterparty-exposure:7}
1.83 million de dollars contre 1.30 million, soit un rapport de 1.41, proche de $\sqrt2 = 1.414$.
\end{solution}

\begin{solution}{exo:rc:counterparty-exposure:8}
Décalages de marge~: en cas de défaut, la banque détient le collatéral du dernier appel de marge satisfait et supporte
le mouvement sur la période de marge en risque (plus le MTA, les litiges et les flux de règlement-livraison)~; la
PFE collatéralisée de l'\omterm{def:rc:counterparty-exposure:netting}{ensemble de compensation} du chapitre reste de 6.45 millions de dollars. La contrepartie
a besoin d'une limite sur cette exposition résiduelle, et sur le risque qu'elle cesse de verser du collatéral.
\end{solution}

\begin{solution}{pb:rc:counterparty-exposure:1}
\textbf{1.} 3.48~\%~; 1.1166 dollars par euro.
\textbf{2.} 4.42 millions de dollars, juste avant que le forward ne se règle à un an.
\textbf{3.} 626\,224 dollars sur un an et 545\,826 sur cinq.
\textbf{4.} Le forward~: 20 millions d'euros à une volatilité de 8~\% présente un risque plus élevé que le swap
durant sa première année, quand peu d'incertitude sur les taux s'est accumulée.
\textbf{5.} Verser des sûretés en espèces chaque jour exige une liquidité et des opérations que la plupart des entreprises n'ont pas
et la comptabilité de couverture privilégie des flux de trésorerie stables.
\textbf{6.} PFE maximale de 3.63 millions de dollars~; EPE de 199\,002 dollars sur un an et de 97\,606 sur cinq.
\textbf{7.} Le forward se règle à un an~: sa valeur quitte l'\omterm{def:rc:counterparty-exposure:netting}{ensemble de compensation} sous forme de paiement
tandis que le collatéral détenu en tient encore compte pendant une période de marge en risque~; si la banque avait
versé du collatéral en contrepartie, celui-ci reste exposé jusqu'à sa restitution.
\textbf{8.} PFE maximale de 4.38 millions de dollars~; EPE sur la durée de vie de 544\,793 dollars.
\textbf{9.} La valeur de l'\omterm{def:rc:counterparty-exposure:netting}{ensemble de compensation} dépasse rarement 5 millions de dollars, si bien que le seuil
est rarement atteint, et presque aucun collatéral ne serait transféré.
\textbf{10.} Des sûretés en espèces ou des titres à verser, une trésorerie capable de faire face aux appels quotidiens et, éventuellement, une
ligne de crédit pour les marges.
\textbf{11.} Non~: sa PFE maximale de 4.42 millions de dollars dépasse la limite.
\textbf{12.} Les nouvelles transactions qui ajoutent de l'exposition au pic à un an exigent une marge sous la limite~; celles
qui la réduisent (le client qui vend des euros à terme) sont bienvenues.
\textbf{13.} Un CSA, un forward plus court, un forward réglé en deux parties, ou une clause de sortie~;
le règlement paiement contre paiement supprime le risque de livraison le jour même.
\textbf{14.} Le forward court sur l'euro de la banque perd environ $20\text{ millions}\times0.11 = \text{USD}~2.2$
millions de dollars~: son exposition à la contrepartie diminue, celle de la contrepartie envers la banque augmentant d'autant.
\textbf{15.} Si l'euro baisse, le forward (la banque vend des euros à 1.1166) rapporte à la banque~: la
banque est exposée à la contrepartie.
\textbf{16.} Les entreprises génèrent l'activité de couverture et disposent de peu d'alternatives~; la banque
facture l'exposition (chapitres~18 à~20) et la limite.
\textbf{17.} Seuls les chiffres collatéralisés changent, d'environ $\sqrt2$ sur l'exposition résiduelle~;
le profil non collatéralisé reste inchangé.
\textbf{18.} Une protection jusqu'au seuil~: l'exposition située sous celui-ci reste non collatéralisée et
doit être évaluée et limitée.
\textbf{19.} Résultat à retenir~: \emph{l'entreprise non collatéralisée}~: PFE maximale de 4.42 millions de dollars et
EPE sur la durée de vie de 545\,826 dollars sans CSA~; avec le CSA à seuil nul, PFE de 3.63 millions de dollars et EPE de 97\,606.
\textbf{20.} Le risque concerne le mouvement de la valeur sur quelques jours, auquel s'ajoutent les flux de règlement-livraison.
\end{solution}

\begin{solution}{iq:rc:counterparty-exposure:1}
EE~: l'exposition positive moyenne à une date, utilisée pour valoriser la CVA. PFE~: un quantile élevé de
l'exposition à une date, utilisé pour fixer les limites de crédit. EPE~: la moyenne temporelle de l'EE, un seul nombre
utilisé dans les mesures réglementaires d'exposition et dans le calcul des fonds propres.
\iqlookfor{les trois définitions et leurs usages.}
\end{solution}

\begin{solution}{iq:rc:counterparty-exposure:2}
Des scénarios communs pour tous les facteurs de risque sur une grille~; réévaluation de chaque transaction sur chaque trajectoire
et date (formules fermées, grilles, régression pour les produits à exercice anticipé), organisée par type de transaction et
vectorisée~; compensation et collatéral par \omterm{def:rc:counterparty-exposure:netting}{ensemble de compensation}~; profils stockés par \omterm{def:rc:counterparty-exposure:netting}{ensemble de compensation} pour
les ajustements. Le temps se dépense dans la réévaluation~; l'agrégation est limitée par la mémoire~; les deux
se parallélisent sur les trajectoires.
\iqlookfor{scénario, réévaluation et agrégation, et où se situe le coût.}
\end{solution}

\begin{solution}{iq:rc:counterparty-exposure:3}
Le délai entre le dernier collatéral reçu et la clôture puis la reconstitution des couvertures de l'\omterm{def:rc:counterparty-exposure:netting}{ensemble
de compensation} défaillant~: constat de l'appel de marge manqué, litiges, périodes de grâce, résiliation et
remplacement. Dix jours ouvrés constituent le minimum réglementaire pour les transactions bilatérales
à marge quotidienne, davantage pour les grands \omterm{def:rc:counterparty-exposure:netting}{ensembles de compensation} ou les ensembles contestés.
\iqlookfor{les étapes de la période et les planchers réglementaires.}
\end{solution}

\begin{solution}{iq:rc:counterparty-exposure:4}
Une banque qui achète une protection sur un État auprès des banques de cet État~; un forward de change dans lequel
une entreprise locale achète des dollars et ferait défaut après une dévaluation. On la mesure en reliant
l'intensité de défaut de la contrepartie aux facteurs de marché de la simulation et en calculant l'exposition
conditionnelle au défaut, ou par des tests de résistance sur des scénarios conjoints.
\iqlookfor{un exemple et une méthode quantitative.}
\end{solution}

\begin{solution}{iq:rc:counterparty-exposure:5}
Sa valeur à chaque date de simulation dépend de la décision d'exercice future~: il faut utiliser une régression
(Longstaff--Schwartz) sur l'état simulé pour estimer la valeur de continuation et l'indicateur d'exercice sur
chaque trajectoire, puis sa valeur sur celle-ci~; après l'exercice, la transaction
devient le swap sous-jacent sur cette trajectoire.
\iqlookfor{réévaluation par régression et exercice dépendant de la trajectoire.}
\end{solution}

\begin{solution}{iq:rc:counterparty-exposure:6}
Elles ont résilié leurs transactions au titre des accords-cadres, les ont compensées, conservé le collatéral et
remplacé leurs couvertures sur un marché sous tension, d'un seul coup~; la plupart des quelque 930\,000 transactions
ont été résiliées en deux mois, certaines étaient encore ouvertes en janvier 2009, et les litiges sur les
montants de résiliation ont duré des années, avec des réclamations supérieures à 50 milliards de dollars.
\iqlookfor{montant de résiliation, coût de remplacement et litiges.}
\end{solution}
